\section{Computational methods}
\label{sec:appendix}

\paragraph{Hardware and software.}
All computations were run on an Apple M4 Max (16 cores, 128\,GB) under macOS, with Python 3.9.6, NumPy 2.0.2, SciPy 1.13.1 and C code compiled with Apple clang 17.
Uniformly random arrangements were generated by Fisher--Yates shuffles (NumPy's default generator; \texttt{xoshiro256**} in C); the Monte Carlo of \Cref{tab:random_energy} uses $10^8$ arrangements per order.
All energies are evaluated from \eqref{eq:Efull}.
The exact statements of \Cref{sec:theory} and the numbers quoted in \Cref{sec:results} are checked in exact integer or rational arithmetic by \texttt{anc/verify\_numbers.py}: the $n=3$ enumeration, all $105{,}888$ single swaps of the $880+8$ magic squares of orders $4$ and $3$ and the swaps of the Siamese squares, the exact means and variances for $n=3,\dots,6$, the null distribution of Loly's invariant, the witnesses, the product-of-inertia values, and the data for \Cref{sec:lowmode_n4}.
The numerical statements of \Cref{sec:balance} are produced by \texttt{anc/balance\_check.py}.

\paragraph{Exact moments and variances.}
The means in \Cref{prop:random_energy_moments} are traces of the covariance matrix $\sigma_n^2G$.
The variances were computed as exact rationals for each $n\le45$ by expanding $\E[E^2]$ over the $15$ equality patterns of a cell quadruple, with the fourth-moment expectations of the uniform permutation expressed through the power sums $p_2,p_4$ of the centred values; the rational functions \eqref{eq:var_low} were then obtained by interpolation separately for odd and even $n$ and verified on $5$--$11$ further orders each, and against the exhaustive $n=3$ values.
Densities of the Gaussian-limit mixtures in \Cref{fig:random_energy} were computed by numerical inversion of the characteristic function, the method of Imhof \cite{imhof1961}.

\paragraph{Meet-in-the-middle enumeration of $L_4$.}
The four statistics $T_1=\sum_{ij}j\,s_{ij}$, $T_2=\sum_{ij}i\,s_{ij}$, $T_3=\sum_i s_{ii}$, $T_4=\sum_i s_{i,n-1-i}$ are additive over cells.
The $16$ cells are split into the top two rows $A$ and the bottom two rows $B$.
For each of the $\binom{16}{8}=12{,}870$ value sets $V$ and each of the $8!$ bijections $V\to A$ the partial vector $(T_1,\dots,T_4)|_A$ is packed into a $27$-bit key; for the complementary set and each bijection $V^c\to B$ the key of (target $-$ partial) is formed; the two key lists are sorted and merged, and the sum over matching keys of the product of multiplicities is $|L_4|$.
This touches $2\times12{,}870\times40{,}320\approx1.04\times10^9$ partial vectors instead of $16!\approx2.1\times10^{13}$ arrangements and runs in $9$\,s on $16$ threads.
Matching pairs are expanded explicitly to obtain the $14{,}639{,}440$ members, which are then classified by row and column sums and by $D_4$-orbit.
Validation is described after \Cref{prop:L4}.

\paragraph{Reproducibility.}
The scripts for every table and figure, the C enumerator and Monte Carlo code, the verification scripts and the JSON files containing the reported numbers are provided as ancillary files in \texttt{anc/}, with a \texttt{README} listing the commands that regenerate them.


\section{Variance closed forms}
\label{sec:app_variance}

For a uniformly random arrangement the variances of the two energies of \Cref{prop:random_energy_moments} are rational
functions of $n$ of the form
\begin{equation}
\begin{aligned}
\Var(\Elow)&=\frac{(n-1)(n^2+1)\,P_{\mathrm o}(n)}{129600\,n^{6}}\ (n\ \text{odd}),&\quad
&\frac{(n^2+1)\,P_{\mathrm e}(n)}{129600\,n^{4}}\ (n\ \text{even});\\
\Var(\Efull)&=\frac{(n-1)(n^2+1)\,Q_{\mathrm o}(n)}{180\,n^{6}}\ (n\ \text{odd}),&\quad
&\frac{(n-2)(n^2+1)\,Q_{\mathrm e}(n)}{180\,n^{5}}\ (n\ \text{even}),
\end{aligned}
\label{eq:var_low}
\end{equation}
with integer polynomials $P_{\mathrm o}$, $P_{\mathrm e}$ of degrees $11$, $10$ and leading coefficient $25$, and
$Q_{\mathrm o}$, $Q_{\mathrm e}$ of degrees $6$, $5$ and leading coefficient $5$, so that $\Var(\Elow)=n^8/5184+O(n^7)$ and
$\Var(\Efull)=n^3/36+O(n^2)$.  The polynomials are listed in the ancillary file \texttt{anc/data\_witnesses\_variances.json}
(and in \texttt{anc/variance\_fits.json}); they were obtained as described in \Cref{sec:appendix}, and at $n=3$ they
give $1532/2187$ and $398/2187$, the exhaustive values, at $n=4$ $14093/1440$ and $527/640$ (\Cref{tab:random_energy}).

\section{A witness for \texorpdfstring{$k^*(5)=4$}{k*(5)=4}}
\label{sec:app_witnesses}

The arrangement
\[
\begin{pmatrix}1&22&11&24&7\\23&3&12&2&25\\18&6&19&17&5\\4&21&14&10&16\\20&9&15&8&13\end{pmatrix}
\]
has all row sums equal to $M(5)=65$ and column sums $66,61,71,61,66$, i.e.\ $\gamma=(1,-4,6,-4,1)=\bigl((-1)^j\binom4j\bigr)_j$;
by \eqref{eq:finite_diff} it satisfies the degree-$3$ conditions of \Cref{def:degree_k} without being semi-magic, so
$k^*(5)=4$.  The analogous witnesses for $6\le n\le13$ (\Cref{rem:kstar_search}) are in the ancillary file
\texttt{anc/data\_witnesses\_variances.json}; \texttt{anc/verify\_numbers.py} verifies them.

